\section{The two characterizations}\label{sec:characterization}

The target page gives two analogues of the Bohr--Mollerup theorem.  They require
care because a first-order arithmetic theory cannot quantify over arbitrary
real functions.  We first prove the ordinary mathematical theorems, then state
their exact coded-function interpretation.

\subsection{Third-derivative positivity for \texorpdfstring{$\log G$}{log G}}

Let \(L(x)=\log G(x)\) for \(x>0\).  From \cref{eq:phi-global},
\begin{align}
 L'''(x)&=2\psi'(x)+(x-1)\psi''(x)\notag\\
 &=2\sum_{k=0}^{\infty}
 \left(\frac1{(x+k)^2}-\frac{x-1}{(x+k)^3}\right)\notag\\
 &=2\sum_{k=0}^{\infty}\frac{k+1}{(x+k)^3}>0.
 \label{eq:Lthird}
\end{align}
For \(x\ge1\), since \(k+1\le x+k\),
\begin{equation}\label{eq:Lthird-decay}
 0<L'''(x)
 \le2\sum_{k=0}^{\infty}\frac1{(x+k)^2}
 \le2\left(\frac1{x^2}+\frac1x\right).
\end{equation}
Thus \(L'''(x)\to0\) with an explicit modulus.

The series proof is constructive.  The trigamma and second-polygamma series
have tails bounded by elementary integrals.  Positivity is witnessed by every
finite partial sum, and \cref{eq:Lthird-decay} is a rational inequality once
\(x\) is replaced by a positive rational lower bound.

\subsection{Vign\'eras uniqueness}

\begin{theorem}[Vign\'eras characterization for Barnes $G$]\label{thm:vigneras}
Let \(f:(0,\infty)\to(0,\infty)\) be three times continuously differentiable
and satisfy
\begin{equation}\label{eq:vigneras-rec}
 f(x+1)=\Gamma(x)f(x),
 \end{equation}
\begin{equation}\label{eq:vigneras-convex}
 \frac{d^3}{dx^3}\log f(x)\ge0\qquad(x\ge1).
\end{equation}
Then \(f(x)=cG(x)\) for a constant \(c>0\).  If \(f(1)=1\), then \(c=1\).
\end{theorem}

\begin{proof}
Put \(g=\log f\) and \(p=g-L\).  The two functional equations imply
\[
 p(x+1)=p(x),
\]
so \(p\) is one-periodic, as are its first three derivatives.  For every fixed
\(x\ge1\) and integer \(n\ge0\),
\[
 0\le g'''(x+n)=L'''(x+n)+p'''(x).
\]
By \cref{eq:Lthird-decay}, \(L'''(x+n)\to0\).  Therefore
\(p'''(x)\ge0\).  Since \(p'''\) is continuous and periodic,
\[
 \int_1^2p'''(x)\dd x=p''(2)-p''(1)=0.
\]
A continuous nonnegative function with zero integral is identically zero, so
\(p'''=0\).  Hence \(p\) is a polynomial of degree at most two.  A periodic
polynomial is constant.  Thus \(g=L+\log c\), and
\(f=cG\).  The normalization at one gives \(c=1\).
\end{proof}

This proof is considerably simpler than a general higher-order convexity
sandwich because the canonical solution has already been constructed and its
third derivative has an explicit decay bound.  It is also particularly well
suited to iPRA.

\begin{proof}[Finite constructive reading]
Fix a coded candidate \(f\), a rational point \(x\), and a precision \(p\).
The code supplies evaluators and continuity moduli for \(g'''\).  Choose
\(n\) so that the rational right side of \cref{eq:Lthird-decay} at \(x+n\) is
at most \(2^{-p}\).  The certified inequality \(g'''(x+n)\ge0\) and periodicity
of \(p'''\) imply
\(p'''(x)\ge-2^{-p}\).  Since this holds for every \(p\), it is the
constructive assertion \(p'''(x)\ge0\).

To infer \(p'''=0\) from zero integral, argue by a finite contradiction.  If a
ball at some rational \(x_0\) has lower endpoint \(\eta>0\), the continuity
modulus supplies a rational interval of positive length on which
\(p'''\ge\eta/2\).  A lower Riemann sum then makes
\(\int_1^2p'''>0\), contradicting the exact periodic endpoint equality.  Thus
every value ball contains zero at every requested precision.  Two certified
integrations show \(p''\), then \(p'\), are constant; periodic endpoint
equalities force those constants to be zero.  Hence \(p\) is constant.
Every step is a finite query about the fixed code.
\end{proof}

\subsection{The corrected asymptotic ratio}

The correct asymptotic law is
\begin{equation}\label{eq:correct-asymp-ratio}
 G(n+x)\sim
 \Gamma(n)^x n^{\binom{x}{2}}G(n)
 \qquad(n\to\infty),
\end{equation}
for each fixed \(x>0\).  We now prove a quantitative version.

\begin{lemma}[Digamma strip estimate]\label{lem:digamma-strip}
For every rational \(X>0\), there is an explicit rational \(C_X\) such that,
for \(n\ge2X+2\) and \(0\le t\le X\),
\begin{equation}\label{eq:digamma-strip}
 \left|
 \psi(n+t)-\log n-\frac{t-1/2}{n}
 \right|\le\frac{C_X}{n^2}.
\end{equation}
\end{lemma}

\begin{proof}
Use the first two terms of the fixed-sector digamma expansion, or differentiate
the Gamma Stirling expansion with its explicit remainder.  Alternatively,
start from
\[
 \psi(y)=\log y-\frac1{2y}
 -\int_0^\infty \frac{2u}{(u^2+y^2)(\e^{2\pi u}-1)}\dd u,
\]
whose last term is bounded by a rational multiple of \(y^{-2}\).  Expand
\(\log(n+t)=\log n+\log(1+t/n)\) to first order with the geometric remainder,
and \((n+t)^{-1}\) about \(n^{-1}\).  On \(t/n\le1/2\), all remainders have
explicit \(C_X/n^2\) bounds.
\end{proof}

Integrate the logarithmic derivative.  For \(0\le x\le X\),
\begin{equation}\label{eq:Ldifference-integral}
 L(n+x)-L(n)=\int_0^x
 \left((n+t-1)\psi(n+t)-(n+t)+c\right)\dd t,
\end{equation}
where \(c=\frac12\log(2\pi)+\frac12\).  Substitute
\cref{eq:digamma-strip} and the one-term Stirling estimate
\begin{equation}\label{eq:loggamma-one-term}
 \log\Gamma(n)=
 \left(n-\frac12\right)\log n-n+\frac12\log(2\pi)
 +O(n^{-1}).
\end{equation}
The main integrand becomes
\[
 \log\Gamma(n)+(t-\tfrac12)\log n,
\]
and the remaining error is bounded by \(D_X/n\) for an explicit rational
\(D_X\).  Integration yields
\begin{equation}\label{eq:ratio-log-error}
 \left|
 \log\frac{G(n+x)}
 {G(n)\Gamma(n)^x n^{x(x-1)/2}}
 \right|
 \le\frac{E_X}{n}.
\end{equation}
Exponentiating gives a relative-error modulus, for example
\begin{equation}\label{eq:ratio-relative-error}
 \left|
 \frac{G(n+x)}{G(n)\Gamma(n)^x n^{x(x-1)/2}}-1
 \right|
 \le \frac{2E_X}{n}
\end{equation}
whenever \(n\ge2E_X\).  This proves \cref{eq:correct-asymp-ratio} uniformly on
every fixed compact \(x\)-interval.

\begin{audit}[Why \(\Gamma(x)^n\) cannot be repaired by convention]
At \(x=1\), the correct formula is exact:
\[
 G(n+1)=\Gamma(n)G(n).
\]
Replacing \(\Gamma(n)^x\) by \(\Gamma(x)^n\) gives \(G(n+1)\sim G(n)\), whose
ratio is \(\Gamma(n)\), not one.  Therefore the source display is a genuine
transposition error, not a branch or indexing convention.
\end{audit}

\subsection{Asymptotic characterization}

\begin{theorem}[Corrected asymptotic characterization]\label{thm:corrected-characterization}
Let \(f:(0,\infty)\to(0,\infty)\) satisfy
\(f(x+1)=\Gamma(x)f(x)\).  Suppose that for every fixed \(x>0\),
\begin{equation}\label{eq:f-asymp}
 f(n+x)\sim\Gamma(n)^x n^{\binom{x}{2}}f(n).
\end{equation}
Then \(f=cG\) for a constant \(c>0\).  Conversely, every \(cG\) satisfies
\cref{eq:f-asymp}.
\end{theorem}

\begin{proof}
The converse is \cref{eq:correct-asymp-ratio}; the constant cancels.  For
uniqueness, put \(h=f/G\).  The functional equations imply
\(h(x+1)=h(x)\).  Divide the asymptotic relation for \(f\) by that for \(G\):
\[
 \frac{h(n+x)}{h(n)}\longrightarrow1.
\]
Periodicity makes the quotient independent of \(n\): it is
\(h(x)/h(1)\).  Therefore \(h(x)=h(1)\) for every \(x>0\), and \(f=cG\).
\end{proof}

The constructive version requires explicit relative-error moduli for the coded
candidate \(f\).  At requested precision \(p\), choose an integer \(n\) large
enough for the \(f\)-modulus and the explicit \(G\)-modulus
\cref{eq:ratio-relative-error}.  Positive lower bounds for the involved values
turn the two relative balls into a ball for \(h(x)/h(1)\).  Periodicity is an
exact finite iteration of the recurrence.  The resulting radius is at most
\(2^{-p}\).  Again there is no object-language quantification over all
functions; there is a theorem schema for each fixed evaluator and modulus code.

\subsection{Comparison of the two characterizations}

The derivative characterization needs a coded \(C^3\) structure and a
nonnegativity checker, but no separate asymptotic input: decay of the canonical
solution forces the periodic discrepancy to be constant.  The asymptotic
characterization needs only positive values and recurrence, but it requires a
relative-error modulus as part of the candidate data.  Both are natural in
FQNF.  The first is preferable for verifying an independently implemented
smooth function; the second is preferable when the implementation already has
an asymptotic evaluator.
